\documentclass[12pt]{amsart}

\usepackage[margin=1.1in]{geometry}
\usepackage[T1]{fontenc}
\usepackage{lmodern}
\usepackage{amsmath,amssymb,amsthm,mathtools}
\usepackage{enumitem}
\usepackage[colorlinks=true,linkcolor=blue,citecolor=blue,urlcolor=blue]{hyperref}
\usepackage{microtype}

\newtheorem{theorem}{Theorem}
\newtheorem{conjecture}{Conjecture}
\newtheorem{proposition}{Proposition}
\newtheorem{lemma}{Lemma}
\newtheorem{corollary}{Corollary}
\newtheorem*{remark}{Remark}
\newtheorem*{example}{Example}

\DeclareMathOperator{\spol}{sp}
\DeclareMathOperator{\num}{num}
\DeclareMathOperator{\den}{den}
\DeclareMathOperator{\rad}{rad}
\newcommand{\floor}[1]{\left\lfloor #1\right\rfloor}
\newcommand{\bbfm}{Ballantine--Beck--Feigon--Maurischat}
\newcommand{\ii}{\mathrm{i}}
\newcommand{\OO}{\mathcal O}
\newcommand{\TT}{\mathcal T}
\newcommand{\BB}{\mathcal B}
\newcommand{\ZZ}{\mathbb Z}

\begin{document}

Let \(\lambda=(\lambda_1,\ldots,\lambda_k)\) be a partition of \(n\).  The
\emph{subsum polynomial} of \(\lambda\) is
\begin{equation}
   \spol(\lambda,x):=\prod_{j=1}^k(1+x^{\lambda_j}).
\end{equation}
Equivalently, if \(m_\lambda(i)\) denotes the multiplicity of the part \(i\) in
\(\lambda\), then we have
\[
   \spol(\lambda,x)=\prod_{i\ge1}(1+x^i)^{m_\lambda(i)}.
\]

Let \(\OO(n)\) be the set of odd partitions of \(n\).  The analogous numerator is
formed from
\[
   h_{O,\lambda}^{(n)}(x)
   :=\prod_{\substack{i\ge1\\ i\text{ odd}}}
      (1+x^i)^{\floor{n/i}-m_\lambda(i)},
\]
with common divisor
\[
   G_O(n,x):=\gcd_{\lambda\in\OO(n)}h_{O,\lambda}^{(n)}(x)
\]
and
\[
   \num_O(n,x):=\frac{1}{G_O(n,x)}
   \sum_{\lambda\in\OO(n)}h_{O,\lambda}^{(n)}(x).
\]

\begin{conjecture}\label{conj:conj8}  Let \(\OO(n)\) be the set of odd partitions of \(n\).  If \(o(m)\) denotes
the largest odd divisor of \(m\), then
\[
   \num_O(n,-1)=o(n!)\qquad(n\ge1).
\]
\end{conjecture}

\end{document}
